\magnification=\magstep1
\hsize=16truecm
\vsize=24truecm
\parindent=0pt
\baselineskip=14pt
\raggedright
\tolerance=10000
\hbadness=10000
\font\tenrm="Liberation Serif:mapping=tex-text" at 10pt
\font\tenbf="Liberation Serif/B:mapping=tex-text" at 10pt
\font\tenit="Liberation Serif/I:mapping=tex-text" at 10pt
\rm
\def\boxnum#1{\hbox{\vrule\vbox{\hrule\hbox{\strut\kern5pt{\bf#1}\kern5pt}\hrule}\vrule}}
\def\problem#1#2{\vskip0pt plus 30pt\penalty-400\bigskip
  \line{\boxnum{#1}\quad #2\hfil}\nobreak\smallskip}
\def\st#1{\medskip\noindent{\bf #1}\enspace\ignorespaces}
\def\col#1#2{\left(\matrix{#1\cr #2\cr}\right)}
\def\mat#1#2#3#4{\left(\matrix{#1&#2\cr #3&#4\cr}\right)}

\centerline{\bf Семинар 13. Линейные системы первого порядка}
\centerline{\bf Собственные значения и фазовая плоскость}
\bigskip

{\bf Памятка: от уравнения к системе.} Для уравнения с постоянными
коэффициентами $y''+py'+qy=0$ положим $x_1=y$, $x_2=y'$. Тогда
$$\cases{x_1'=x_2,\cr x_2'=-qx_1-px_2,\cr}\qquad
X'=AX,\quad X=\col{x_1}{x_2},\quad A=\mat{0}{1}{-q}{-p}.$$
Начальные условия $y(0)=a$, $y'(0)=b$ становятся $X(0)=\col ab$.
Штрих далее означает производную по времени $t$.

\st{Экспоненциальное решение.} Для системы $X'=AX$ с постоянной
действительной матрицей ищем $X=e^{\lambda t}v$, где $v\ne0$ -- постоянный
столбец. Подстановка даёт $\lambda e^{\lambda t}v=e^{\lambda t}Av$,
поэтому
$$Av=\lambda v,\qquad (A-\lambda I)v=0,\qquad \det(A-\lambda I)=0.$$
Сначала находим собственные значения $\lambda$, затем для каждого --
ненулевой собственный вектор $v$.

\st{Два различных действительных значения.} Для соответствующих векторов
$v_1,v_2$ общее действительное решение имеет вид
$$X=C_1e^{\lambda_1t}v_1+C_2e^{\lambda_2t}v_2.$$
Подстановка $t=0$ даёт систему на $C_1,C_2$. Если оба значения отрицательны,
траектории стремятся к началу координат (устойчивый узел); если оба
положительны, удаляются (неустойчивый узел). При разных знаках получается
седло: вдоль направления отрицательного значения движение к началу,
вдоль направления положительного -- от него.

\st{Комплексная пара.} Пусть $\lambda=\alpha+i\beta$, $\beta>0$,
и $v=a+ib$, где $a,b$ -- действительные столбцы. Тогда действительная
и мнимая части комплексного решения дают два действительных решения:
$$X_1=e^{\alpha t}(a\cos\beta t-b\sin\beta t),\qquad
X_2=e^{\alpha t}(a\sin\beta t+b\cos\beta t).$$
Общее решение $X=C_1X_1+C_2X_2$. При $\alpha<0$ траектории
закручиваются к началу (устойчивый фокус), при $\alpha>0$ -- от него;
при $\alpha=0$ получаются замкнутые траектории (центр).

\st{Фазовый эскиз.} Оси -- компоненты $x$ и $y$ системы; точка $(x,y)$
движется со скоростью $(x',y')$. Отметьте начало координат и направления
собственных векторов, если они действительны. Стрелки на траекториях
покажите при возрастании $t$; для вращения найдите вектор скорости в
удобной точке, например $(1,0)$. Проверяйте найденное решение подстановкой
в оба уравнения и в начальные условия.

\vfil\eject
\centerline{\bf Классная работа: пример преподавателя}
\bigskip
\problem{1}{$x'=x+2y,\quad y'=2x+y,\quad x(0)=2,\ y(0)=0$}
\st{Шаг 1. Матрица и значения.} Записываем
$$X'=AX,\quad X=\col xy,\quad A=\mat1221.$$
$$\det(A-\lambda I)=\left|\matrix{1-\lambda&2\cr 2&1-\lambda\cr}\right|
=(1-\lambda)^2-4=(\lambda-3)(\lambda+1)=0.$$
Значит, $\lambda_1=3$, $\lambda_2=-1$.

\st{Шаг 2. Собственные векторы.} Для $\lambda_1=3$ система
$(A-3I)v=0$ даёт $-2v_x+2v_y=0$, поэтому можно взять
$v_1=\col11$. Для $\lambda_2=-1$ получаем $2v_x+2v_y=0$,
поэтому $v_2=\col1{-1}$.

\st{Шаг 3. Общее решение и начальные условия.}
$$X=C_1e^{3t}\col11+C_2e^{-t}\col1{-1},$$
$$x=C_1e^{3t}+C_2e^{-t},\qquad y=C_1e^{3t}-C_2e^{-t}.$$
При $t=0$ имеем $C_1+C_2=2$, $C_1-C_2=0$; следовательно,
$C_1=C_2=1$ и
$$x=e^{3t}+e^{-t},\qquad y=e^{3t}-e^{-t}.$$
\st{Проверка.} Производные равны $x'=3e^{3t}-e^{-t}$,
$y'=3e^{3t}+e^{-t}$. Подстановка даёт
$$x+2y=3e^{3t}-e^{-t}=x',\qquad
2x+y=3e^{3t}+e^{-t}=y'.$$
Кроме того, $x(0)=1+1=2$, $y(0)=1-1=0$.

\st{Фазовый эскиз.} Значения разных знаков: начало координат -- седло.
На прямой $y=x$ стрелки направлены от начала, на прямой $y=-x$ -- к нему.
Траектория из $(2,0)$ в начальный момент имеет скорость $(2,4)$;
отметьте её стрелкой и продолжите с учётом направлений седла.

\vfil\eject
\centerline{\bf Классная работа: у доски}
\bigskip
Решите задачи самостоятельно: найдите общее решение, затем решение с
начальными условиями; проверьте его подстановкой. Для каждой системы
сделайте фазовый эскиз и укажите направление движения.
\problem{2}{$x'=2x+y,\quad y'=x+2y,\quad x(0)=1,\ y(0)=0$}
\problem{3}{$y''+4y=0,\quad y(0)=1,\ y'(0)=0$}
В задаче 3 сначала перейдите к системе для $x_1=y$, $x_2=y'$;
нарисуйте фазовый портрет в координатах $(y,y')$.
\problem{4}{$x'=-x-2y,\quad y'=2x-y,\quad x(0)=1,\ y(0)=0$}

\bigskip
\centerline{\bf Домашняя практика}
\bigskip
Найдите общее действительное решение каждой системы; в задачах с начальными
условиями найдите также соответствующее решение. Проверьте подстановкой.
\problem{1}{$x'=3x+y,\quad y'=x+3y,\quad x(0)=1,\ y(0)=1$}
\problem{2}{$x'=-x+y,\quad y'=2y$}
Для задачи 2 сделайте фазовый эскиз со стрелками.
\problem{3}{$x'=-2x-y,\quad y'=x-2y,\quad x(0)=1,\ y(0)=0$}
\problem{4}{$y''+3y'+2y=0$}
В задаче 4 сначала запишите эквивалентную систему для $(y,y')$ и решите её.

\medskip
В моделях «хищник -- жертва» также изучают движение на фазовой плоскости;
там система обычно нелинейна, но направление скоростей помогает строить
эскиз и понимать поведение популяций.
\bye
